\subsection[Challenge 32: Rayleigh-Darcy convection with mixed boundary conditions]{Challenge 32: Rayleigh-Darcy convection with mixed boundary conditions}
\textbf{Problem ID:} \texttt{Challenge\_32\_main}\\
\textbf{Primary inferred discipline:} Fluid Dynamics\\
\textbf{Python implementation:} \texttt{challenges/32/answer.py}

\subsubsection*{Problem}
\paragraph*{Problem setup}
Consider natural convection in a porous medium modeled by Darcy's law (i.e., Rayleigh-Darcy convection) with horizontal periodic boundary conditions. The boundary condition on the top wall is free (constant pressure) and at constant temperature. The boundary condition on the bottom wall is impermeable and has constant heat flux.\textcolor{blue}{The flow follows the Boussinesq approxiamtion: the density varies with temperature only in buoyancy term, and is taken asa a constant elsewhere.}

\paragraph*{Main problem}

What is the critical Rayleigh number (allowing $\pm 0.05$ error) and associated critical horizontal wavenumber (allowing $\pm 0.02$ error), where, above this critical Rayleigh number, the conduction base state is linearly unstable? Moreover, consider eigenfunction of vertical velocity w(z) and temperature T(z) associated with this critical Rayleigh number and crical horizontal wavenumber, what is the value of w(z)/T(z) at z=0.67365, where the bottom wall is at z=0 and the top wall is at z=1?
